Capítulo  4
 Armadilhas  e  chamadas  de  sistema
 Existem  três  tipos  de  eventos  que  fazem  com  que  a  CPU  interrompa  a  execução  normal  de  instruções  e  force  a  
transferência  de  controle  para  um  código  especial  que  lida  com  o  evento.  Uma  situação  é  uma  chamada  de  sistema,  
quando  um  programa  de  usuário  executa  a  instrução  ecall  para  solicitar  que  o  kernel  faça  algo  por  ele.  Outra  situação  é  
uma  exceção:  uma  instrução  (de  usuário  ou  kernel)  faz  algo  ilegal,  como  dividir  por  zero  ou  usar  um  endereço  virtual  
inválido.  A  terceira  situação  é  uma  interrupção  de  dispositivo,  quando  um  dispositivo  sinaliza  que  precisa  de  atenção,  por  
exemplo,  quando  o  hardware  do  disco  conclui  uma  solicitação  de  leitura  ou  gravação.
 Este  livro  usa  "trap"  como  um  termo  genérico  para  essas  situações.  Normalmente,  qualquer  código  que  estivesse  
em  execução  no  momento  da  trap  precisará  ser  retomado  posteriormente  e  não  precisaria  estar  ciente  de  que  algo  
especial  aconteceu.  Ou  seja,  frequentemente  queremos  que  as  traps  sejam  transparentes;  isso  é  particularmente  
importante  para  interrupções  de  dispositivo,  que  o  código  interrompido  normalmente  não  espera.  A  sequência  usual  é  
que  uma  trap  força  uma  transferência  de  controle  para  o  kernel;  o  kernel  salva  registradores  e  outros  estados  para  que  a  
execução  possa  ser  retomada;  o  kernel  executa  o  código  manipulador  apropriado  (por  exemplo,  uma  implementação  de  
chamada  de  sistema  ou  driver  de  dispositivo);  o  kernel  restaura  o  estado  salvo  e  retorna  da  trap;  e  o  código  original  
continua  de  onde  parou.
 O  Xv6  lida  com  todas  as  armadilhas  no  kernel;  as  armadilhas  não  são  entregues  ao  código  do  usuário.  Lidar  com  
armadilhas  no  kernel  é  natural  para  chamadas  de  sistema.  Faz  sentido  para  interrupções,  já  que  o  isolamento  exige  que  
apenas  o  kernel  tenha  permissão  para  usar  dispositivos  e  porque  o  kernel  é  um  mecanismo  conveniente  para  compartilhar  
dispositivos  entre  vários  processos.  Também  faz  sentido  para  exceções,  já  que  o  Xv6  responde  a  todas  as  exceções  do  
espaço  do  usuário  eliminando  o  programa  ofensivo.
 O  tratamento  de  traps  Xv6  ocorre  em  quatro  etapas:  ações  de  hardware  executadas  pela  CPU  RISC-V,  algumas  
instruções  de  montagem  que  preparam  o  caminho  para  o  código  C  do  kernel,  uma  função  C  que  decide  o  que  fazer  com  
o  trap  e  a  chamada  de  sistema  ou  rotina  de  serviço  do  driver  de  dispositivo.  Embora  a  semelhança  entre  os  três  tipos  de  
traps  sugira  que  um  kernel  poderia  lidar  com  todos  os  traps  com  um  único  caminho  de  código,  torna-se  conveniente  ter  
código  separado  para  três  casos  distintos:  traps  do  espaço  do  usuário,  traps  do  espaço  do  kernel  e  interrupções  de  
temporizador.  O  código  do  kernel  (assembler  ou  C)  que  processa  um  trap  é  frequentemente  chamado  de  manipulador;  
as  primeiras  instruções  do  manipulador  são  geralmente  escritas  em  assembler  (em  vez  de  C)  e,  às  vezes,  são  chamadas  
de  vetor.
 43
Machine Translated by Google
 4.1  Máquinas  de  armadilha  RISC-V
 Cada  CPU  RISC-V  possui  um  conjunto  de  registradores  de  controle  que  o  kernel  escreve  para  informar  
à  CPU  como  lidar  com  traps,  e  que  o  kernel  pode  ler  para  descobrir  se  um  trap  ocorreu.  Os  
documentos  RISC-V  contêm  a  história  completa  [3].  riscv.h  (kernel/riscv.h:1)  contém  definições  que  o  xv6  usa.
 Aqui  está  um  esboço  dos  registros  mais  importantes:
 •  stvec:  O  kernel  escreve  o  endereço  do  seu  manipulador  de  trap  aqui;  o  RISC-V  salta  para  o
 endereço  em  stvec  para  manipular  uma  armadilha.
 •
 sepc:  Quando  ocorre  uma  trap,  o  RISC-V  salva  o  contador  de  programa  aqui  (já  que  o  pc  é  então  sobrescrito  pelo  
valor  em  stvec).  A  instrução  sret  (retorno  da  trap)  copia  o  sepc  para  o  pc.  O  kernel  pode  escrever  sepc  para  controlar  
para  onde  o  sret  vai.
 •  causa:  o  RISC-V  coloca  um  número  aqui  que  descreve  o  motivo  da  armadilha.
 •  sscratch:  O  código  do  manipulador  de  traps  usa  sscratch  para  ajudar  a  evitar  a  substituição  de  registros  de  usuários
 antes  de  salvá-los.
 •  sstatus:  O  bit  SIE  em  sstatus  controla  se  as  interrupções  do  dispositivo  estão  habilitadas.  Se  o  kernel  limpar  o  SIE,  
o  RISC-V  adiará  as  interrupções  do  dispositivo  até  que  o  kernel  configure  o  SIE.  O  bit  SPP  indica  se  uma  
interceptação  veio  do  modo  usuário  ou  do  modo  supervisor  e  controla  para  qual  modo  sret  retorna.
 Os  registradores  acima  referem-se  a  traps  manipulados  no  modo  supervisor  e  não  podem  ser  lidos  ou  escritos  no  modo  
usuário.  Existe  um  conjunto  semelhante  de  registradores  de  controle  para  traps  manipulados  no  modo  máquina;  o  xv6  os  
utiliza  apenas  para  o  caso  especial  de  interrupções  de  temporizador.
 Cada  CPU  em  um  chip  multinúcleo  tem  seu  próprio  conjunto  desses  registradores,  e  mais  de  uma  CPU  pode  estar  
manipulando  uma  armadilha  a  qualquer  momento.
 Quando  precisa  forçar  uma  armadilha,  o  hardware  RISC-V  faz  o  seguinte  para  todos  os  tipos  de  armadilha  (exceto  
interrupções  de  temporizador):
 1.  Se  a  armadilha  for  uma  interrupção  de  dispositivo  e  o  bit  SIE  sstatus  estiver  limpo,  não  faça  nada
 seguindo.
 2.  Desabilite  as  interrupções  limpando  o  bit  SIE  em  sstatus.
 3.  Copie  o  pc  para  sepc.
 4.  Salve  o  modo  atual  (usuário  ou  supervisor)  no  bit  SPP  em  sstatus.
 5.  Defina  scause  para  refletir  a  causa  da  armadilha.
 6.  Defina  o  modo  como  supervisor.
 7.  Copie  stvec  para  o  pc.
 44
Machine Translated by Google
 8.  Comece  a  executar  no  novo  PC.
 Observe  que  a  CPU  não  alterna  para  a  tabela  de  páginas  do  kernel,  não  alterna  para  uma  pilha  no  kernel  e  não  salva  
nenhum  registrador  além  do  pc.  O  software  do  kernel  deve  executar  essas  tarefas.
 Um  motivo  pelo  qual  a  CPU  realiza  um  trabalho  mínimo  durante  uma  armadilha  é  para  fornecer  flexibilidade  ao  software;  por  
exemplo,  alguns  sistemas  operacionais  omitem  uma  troca  de  tabela  de  páginas  em  algumas  situações  para  aumentar  o  
desempenho  da  armadilha.
 Vale  a  pena  considerar  se  algum  dos  passos  listados  acima  poderia  ser  omitido,  talvez  em  busca  de  armadilhas  mais  
rápidas.  Embora  existam  situações  em  que  uma  sequência  mais  simples  possa  funcionar,  muitas  das  etapas  seriam  perigosas  
se  omitidas  em  geral.  Por  exemplo,  suponha  que  a  CPU  não  alterne  os  contadores  de  programa.  Então,  uma  armadilha  do  
espaço  do  usuário  poderia  alternar  para  o  modo  supervisor  enquanto  ainda  executa  instruções  do  usuário.  Essas  instruções  do  
usuário  poderiam  quebrar  o  isolamento  usuário/kernel,  por  exemplo,  modificando  o  registrador  satp  para  apontar  para  uma  
tabela  de  páginas  que  permitisse  o  acesso  a  toda  a  memória  física.
 Portanto,  é  importante  que  a  CPU  alterne  para  um  endereço  de  instrução  especificado  pelo  kernel,  ou  seja,  stvec.
 4.2  Armadilhas  do  espaço  do  usuário
 O  Xv6  lida  com  armadilhas  de  forma  diferente,  dependendo  se  elas  ocorrem  durante  a  execução  no  kernel  ou  no  código  
do  usuário.  Aqui  está  a  história  das  armadilhas  do  código  do  usuário;  a  Seção  4.5  descreve  as  armadilhas  do  código  do  
kernel.
 Uma  armadilha  pode  ocorrer  durante  a  execução  no  espaço  do  usuário  se  o  programa  do  usuário  fizer  uma  chamada  de  
sistema  (instrução  ecall),  ou  algo  ilegal,  ou  se  um  dispositivo  interromper.  O  caminho  de  alto  nível  de  uma  armadilha  do  espaço  
do  usuário  é  uservec  (kernel/trampoline.S:21).  então  usertrap  (kernel/trap.c:37);  e  ao  retornar,  usertrapret  (kernel/trap.c:90)  e  
então  userret  (kernel/trampoline.S:101).
 Uma  grande  limitação  no  design  do  tratamento  de  traps  do  xv6  é  o  fato  de  que  o  hardware  RISC-V  não  alterna  entre  
tabelas  de  páginas  ao  forçar  um  trap.  Isso  significa  que  o  endereço  do  manipulador  de  traps  em  stvec  deve  ter  um  mapeamento  
válido  na  tabela  de  páginas  do  usuário,  já  que  essa  é  a  tabela  de  páginas  em  vigor  quando  o  código  de  tratamento  de  traps  
começa  a  ser  executado.  Além  disso,  o  código  de  tratamento  de  traps  do  xv6  precisa  alternar  para  a  tabela  de  páginas  do  
kernel;  para  poder  continuar  a  execução  após  essa  troca,  a  tabela  de  páginas  do  kernel  também  deve  ter  um  mapeamento  para  
o  manipulador  apontado  por  stvec.
 O  Xv6  atende  a  esses  requisitos  usando  uma  página  de  trampolim.  A  página  de  trampolim  contém  uservec,
 O  código  de  manipulação  de  traps  xv6  para  o  qual  stvec  aponta.  A  página  do  trampolim  é  mapeada  na  tabela  de  páginas  de  
cada  processo  no  endereço  TRAMPOLINE,  que  está  no  topo  do  espaço  de  endereço  virtual,  de  modo  que  estará  acima  da  
memória  que  os  programas  usam  para  si  mesmos.  A  página  do  trampolim  também  é  mapeada  no  endereço  TRAMPOLINE  na  
tabela  de  páginas  do  kernel.  Veja  as  Figuras  2.3  e  3.3.  Como  a  página  do  trampolim  é  mapeada  na  tabela  de  páginas  do  
usuário,  sem  o  sinalizador  PTE_U,  as  traps  podem  começar  a  ser  executadas  lá  no  modo  supervisor.  Como  a  página  do  
trampolim  é  mapeada  no  mesmo  endereço  no  espaço  de  endereço  do  kernel,  o  manipulador  de  traps  pode  continuar  a  execução  
após  alternar  para  a  tabela  de  páginas  do  kernel.
 O  código  para  o  manipulador  de  armadilhas  uservec  está  em  trampoline.S  (kernel/trampoline.S:21).  Quando  o  uservec  
inicia,  todos  os  32  registradores  contêm  valores  pertencentes  ao  código  de  usuário  interrompido.  Esses  32  valores  precisam  
ser  salvos  em  algum  lugar  na  memória,  para  que  possam  ser  restaurados  quando  a  armadilha  retornar  ao  espaço  do  usuário.  
Armazenar  na  memória  requer  o  uso  de  um  registrador  para  armazenar  o  endereço,  mas  neste  ponto  não  há
 45
Machine Translated by Google
 Não  há  registradores  de  uso  geral  disponíveis!  Felizmente,  o  RISC-V  oferece  uma  ajuda  na  forma  do  registrador  
sscratch.  A  instrução  csrw  no  início  de  uservec  salva  a0  em  sscratch.
 Agora  o  uservec  tem  um  registrador  (a0)  para  brincar.
 A  próxima  tarefa  do  uservec  é  salvar  os  32  registradores  de  usuário.  O  kernel  aloca,  para  cada  processo,  uma  
página  de  memória  para  uma  estrutura  trapframe  que  (entre  outras  coisas)  tem  espaço  para  salvar  os  32  registradores  
de  usuário  (kernel/proc.h:43).  Como  satp  ainda  se  refere  à  tabela  de  páginas  do  usuário,  uservec  precisa  que  o  
trapframe  seja  mapeado  no  espaço  de  endereço  do  usuário.  Xv6  mapeia  o  trapframe  de  cada  processo  no  endereço  
virtual  TRAPFRAME  na  tabela  de  páginas  do  usuário  desse  processo;  TRAPFRAME  está  logo  abaixo  de  TRAMPOLINE.
 O  p->trapframe  do  processo  também  aponta  para  o  trapframe,  embora  em  seu  endereço  físico,  para  que  o  kernel  
possa  usá-lo  por  meio  da  tabela  de  páginas  do  kernel.
 Assim,  uservec  carrega  o  endereço  TRAPFRAME  em  a0  e  salva  todos  os  registros  do  usuário  lá,  incluindo  o  a0  
do  usuário,  lidos  do  zero.
 O  trapframe  contém  o  endereço  da  pilha  do  kernel  do  processo  atual,  o  hartid  da  CPU  atual,  o  endereço  da  função  
usertrap  e  o  endereço  da  tabela  de  páginas  do  kernel.  uservec  recupera  esses  valores,  alterna  satp  para  a  tabela  de  
páginas  do  kernel  e  chama  usertrap.
 O  trabalho  do  usertrap  é  determinar  a  causa  da  armadilha,  processá-la  e  retornar  (kernel/-  trap.c:37).  Primeiro,  ele  
altera  stvec  para  que  uma  trap  enquanto  estiver  no  kernel  seja  manipulada  por  kernelvec  em  vez  de  uservec.  Ele  salva  
o  registrador  sepc  (o  contador  de  programa  do  usuário  salvo),  porque  usertrap  pode  chamar  yield  para  alternar  para  a  
thread  do  kernel  de  outro  processo,  e  esse  processo  pode  retornar  ao  espaço  do  usuário,  no  processo  do  qual  ele  
modificará  sepc.  Se  a  trap  for  uma  chamada  de  sistema,  usertrap  chama  syscall  para  manipulá-la;  se  for  uma  
interrupção  de  dispositivo,  devintr;  caso  contrário,  é  uma  ex-ception,  e  o  kernel  mata  o  processo  com  falha.  O  caminho  
da  chamada  de  sistema  adiciona  quatro  ao  contador  de  programa  do  usuário  salvo  porque  o  RISC-V,  no  caso  de  uma  
chamada  de  sistema,  deixa  o  ponteiro  do  programa  apontando  para  a  instrução  ecall,  mas  o  código  do  usuário  precisa  
retomar  a  execução  na  instrução  subsequente.
 Na  saída,  o  usertrap  verifica  se  o  processo  foi  encerrado  ou  se  deve  render  a  CPU  (se  essa  armadilha  for  uma  
interrupção  de  timer).
 O  primeiro  passo  para  retornar  ao  espaço  do  usuário  é  a  chamada  para  usertrapret  (kernel/trap.c:90).  Esta  
função  configura  os  registradores  de  controle  RISC-V  para  se  prepararem  para  uma  futura  captura  a  partir  do  
espaço  do  usuário.  Isso  envolve  alterar  stvec  para  se  referir  a  uservec,  preparar  os  campos  do  trapframe  dos  
quais  uservec  depende  e  definir  sepc  para  o  contador  de  programa  do  usuário  salvo  anteriormente.  No  final,  
usertrapret  chama  userret  na  página  do  trampolim  mapeada  nas  tabelas  de  páginas  do  usuário  e  do  kernel;  o  
motivo  é  que  o  código  assembly  em  userret  alternará  as  tabelas  de  páginas.
 A  chamada  de  usertrapret  para  userret  passa  um  ponteiro  para  a  tabela  de  páginas  de  usuário  do  processo  em  
a0  (kernel/trampoline.S:101).  userret  alterna  o  satp  para  a  tabela  de  páginas  do  usuário  do  processo.  Lembre-se  de  
que  a  tabela  de  páginas  do  usuário  mapeia  tanto  a  página  do  trampolim  quanto  a  TRAPFRAME,  mas  nada  mais  do  kernel.
 O  mapeamento  da  página  do  trampolim  no  mesmo  endereço  virtual  nas  tabelas  de  páginas  do  usuário  e  do  kernel  
permite  que  o  userret  continue  executando  após  a  alteração  do  satp.  A  partir  deste  ponto,  os  únicos  dados  que  o  
userret  pode  usar  são  o  conteúdo  do  registrador  e  o  conteúdo  do  trapframe.  O  userret  carrega  o  endereço  TRAPFRAME  
em  a0,  restaura  os  registradores  de  usuário  salvos  do  trapframe  via  a0,  restaura  o  usuário  a0  salvo  e  executa  sret  
para  retornar  ao  espaço  do  usuário.
 46
Machine Translated by Google
 4.3  Código:  Chamando  chamadas  de  sistema
 O  capítulo  2  terminou  com  initcode.S  invocando  a  chamada  de  sistema  exec  (user/initcode.S:11).  Vamos  ver  como  a  chamada  
do  usuário  chega  à  implementação  da  chamada  do  sistema  exec  no  kernel.
 initcode.S  coloca  os  argumentos  para  exec  nos  registradores  a0  e  a1,  e  coloca  o  número  da  chamada  do  sistema  em  a7.  
Os  números  da  chamada  do  sistema  correspondem  às  entradas  no  array  syscalls,  uma  tabela  de  ponteiros  de  função  (kernel/
 syscall.c:107).  A  instrução  ecall  é  capturada  no  kernel  e  faz  com  que  uservec,  usertrap  e  então  syscall  sejam  executados,  
como  vimos  acima.
 chamada  de  sistema  (kernel/syscall.c:132)  recupera  o  número  de  chamada  do  sistema  do  a7  salvo  no  trapframe
 e  o  utiliza  para  indexar  chamadas  de  sistema.  Para  a  primeira  chamada  de  sistema,  a7  contém  SYS_exec  (kernel /syscall.h:8),  
resultando  em  uma  chamada  para  a  função  de  implementação  de  chamada  do  sistema  sys_exec.
 Quando  sys_exec  retorna,  syscall  registra  seu  valor  de  retorno  em  p->trapframe->a0.  Isso  fará  com  que  a  chamada  
original  do  espaço  do  usuário  para  exec()  retorne  esse  valor,  já  que  a  convenção  de  chamada  C  no  RISC-V  coloca  os  valores  
de  retorno  em  a0.  Chamadas  de  sistema  convencionalmente  retornam  números  negativos  para  indicar  erros  e  zero  ou  números  
positivos  para  sucesso.  Se  o  número  da  chamada  de  sistema  for  inválido,  syscall  imprime  um  erro  e  retorna  -1.
 4.4  Código:  Argumentos  de  chamada  de  sistema
 Implementações  de  chamadas  de  sistema  no  kernel  precisam  encontrar  os  argumentos  passados  pelo  código  do  usuário.  
Como  o  código  do  usuário  chama  funções  wrapper  de  chamadas  de  sistema,  os  argumentos  estão  inicialmente  onde  a  
convenção  de  chamada  RISC-VC  os  coloca:  em  registradores.  O  código  de  trap  do  kernel  salva  os  registradores  do  
usuário  no  quadro  de  trap  do  processo  atual,  onde  o  código  do  kernel  pode  encontrá-los.  As  funções  do  kernel  argint,  
argaddr  e  argfd  recuperam  o  n  'ésimo  argumento  de  chamada  de  sistema  do  quadro  de  trap  como  um  inteiro,  um  ponteiro  
ou  um  descritor  de  arquivo.  Todas  elas  chamam  argraw  para  recuperar  o  registrador  de  usuário  salvo  apropriado  (kernel/syscall.c:34).
 Algumas  chamadas  de  sistema  passam  ponteiros  como  argumentos,  e  o  kernel  deve  usar  esses  ponteiros  para  ler  ou  
escrever  na  memória  do  usuário.  A  chamada  de  sistema  exec,  por  exemplo,  passa  ao  kernel  um  conjunto  de  ponteiros  
referentes  a  argumentos  de  string  no  espaço  do  usuário.  Esses  ponteiros  apresentam  dois  desafios.  Primeiro,  o  programa  do  
usuário  pode  estar  com  bugs  ou  ser  malicioso,  e  pode  passar  ao  kernel  um  ponteiro  inválido  ou  um  ponteiro  com  a  intenção  
de  induzir  o  kernel  a  acessar  a  memória  do  kernel  em  vez  da  memória  do  usuário.  Segundo,  os  mapeamentos  de  tabelas  de  
páginas  do  kernel  xv6  não  são  os  mesmos  que  os  mapeamentos  de  tabelas  de  páginas  do  usuário,  portanto,  o  kernel  não  
pode  usar  instruções  comuns  para  carregar  ou  armazenar  a  partir  de  endereços  fornecidos  pelo  usuário.
 O  kernel  implementa  funções  que  transferem  dados  com  segurança  de  e  para  endereços  fornecidos  pelo  usuário.  fetchstr  
é  um  exemplo  (kernel/syscall.c:25).  Chamadas  de  sistema  de  arquivos  como  exec  usam  fetchstr  para  recuperar  argumentos  
de  nome  de  arquivo  de  string  do  espaço  do  usuário.  fetchstr  chama  copyinstr  para  fazer  o  trabalho  pesado.  copyinstr  (kernel/
 vm.c:403)  Copia  até  max  bytes  para  dst  do  endereço  virtual  srcva  na  tabela  de  páginas  do  usuário  pagetable.  Como  pagetable  
não  é  a  tabela  de  páginas  atual,  copyinstr  usa  walkaddr  (que  chama  walk)  para  procurar  srcva  na  pagetable,  retornando  o  
endereço  físico  pa0.  O  kernel  mapeia  cada  endereço  físico  de  RAM  para  o  endereço  virtual  do  kernel  correspondente,  de  
modo  que  copyinstr  pode  copiar  diretamente  bytes  de  string  de  pa0  para  dst.  walkaddr  (kernel/vm.c:109)  verifica  se  o  endereço  
virtual  fornecido  pelo  usuário  faz  parte  do  espaço  de  endereço  do  usuário  do  processo,  para  que  os  programas
 47
Machine Translated by Google
 não  consegue  enganar  o  kernel  para  que  ele  leia  outra  memória.  Uma  função  semelhante,  copyout,  copia  dados  do  kernel  para  
um  endereço  fornecido  pelo  usuário.
 4.5  Armadilhas  do  espaço  do  kernel
 O  Xv6  configura  os  registradores  de  trap  da  CPU  de  forma  um  pouco  diferente,  dependendo  se  o  código  do  usuário  ou  do  kernel  
está  em  execução.  Quando  o  kernel  está  em  execução  em  uma  CPU,  o  kernel  aponta  stvec  para  o  código  assembly  em  kernelvec  
(kernel/kernelvec.S:12).  Como  o  xv6  já  está  no  kernel,  o  kernelvec  pode  contar  com  satp  sendo  definido  na  tabela  de  páginas  do  
kernel  e  com  o  ponteiro  da  pilha  se  referindo  a  uma  pilha  do  kernel  válida.  O  kernelvec  envia  todos  os  32  registradores  para  a  
pilha,  de  onde  ele  os  restaurará  posteriormente  para  que  o  código  do  kernel  interrompido  possa  ser  retomado  sem  perturbações.
 kernelvec  salva  os  registradores  na  pilha  da  thread  do  kernel  interrompida,  o  que  faz  sentido  porque  os  valores  dos  
registradores  pertencem  a  essa  thread.  Isso  é  particularmente  importante  se  a  armadilha  causar  uma  troca  para  uma  thread  
diferente  –  nesse  caso,  a  armadilha  retornará  da  pilha  da  nova  thread,  deixando  os  registradores  salvos  da  thread  interrompida  
em  segurança  em  sua  pilha.
 kernelvec  salta  para  kerneltrap  (kernel/trap.c:135)  Após  salvar  os  registradores,  o  kerneltrap  está  
preparado  para  dois  tipos  de  armadilhas:  interrupções  de  dispositivo  e  exceções.  Ele  chama  devintr  (kernel/  -trap.c:178)  para  verificar  e  lidar  com  o  primeiro.  Se  a  armadilha  não  for  uma  interrupção  de  dispositivo,  deve  
ser  uma  exceção,  e  isso  é  sempre  um  erro  fatal  se  ocorrer  no  kernel  xv6;  o  kernel  dispara  pânico  e  interrompe  
a  execução.
 Se  o  kerneltrap  foi  chamado  devido  a  uma  interrupção  do  temporizador  e  uma  thread  do  kernel  de  um  processo  está  em  
execução  (em  oposição  a  uma  thread  do  escalonador),  o  kerneltrap  chama  o  método  yield  para  dar  a  outras  threads  a  chance  
de  executar.  Em  algum  momento,  uma  dessas  threads  irá  executar  o  método  yield,  permitindo  que  nossa  thread  e  sua  kerneltrap  
sejam  retomadas.  O  Capítulo  7  explica  o  que  acontece  no  método  yield.
 Quando  o  trabalho  do  kerneltrap  termina,  ele  precisa  retornar  ao  código  que  foi  interrompido  pela  armadilha.  Como  um  yield  
pode  ter  perturbado  o  sepc  e  o  modo  anterior  em  sstatus,  o  kerneltrap  os  salva  ao  iniciar.  Agora,  ele  restaura  esses  registradores  
de  controle  e  retorna  ao  kernelvec  (kernel/kernelvec.S:50).  kernelvec  retira  os  registradores  salvos  da  pilha  e  executa  sret,  que  
copia  sepc  para  pc  e  retoma  o  código  do  kernel  interrompido.
 Vale  a  pena  pensar  em  como  o  retorno  do  trap  acontece  se  o  kerneltrap  chamar  yield  devido  a  uma  interrupção  do  
temporizador.
 O  Xv6  define  o  stvec  de  uma  CPU  como  kernelvec  quando  essa  CPU  entra  no  kernel  pelo  espaço  do  
usuário;  você  pode  ver  isso  em  usertrap  (kernel/trap.c:29).  Há  uma  janela  de  tempo  em  que  o  kernel  inicia  a  
execução,  mas  stvec  ainda  está  definido  como  uservec,  e  é  crucial  que  nenhuma  interrupção  de  dispositivo  
ocorra  durante  essa  janela.  Felizmente,  o  RISC-V  sempre  desabilita  as  interrupções  quando  começa  a  receber  
uma  captura,  e  o  xv6  não  as  habilita  novamente  até  que  stvec  seja  definido.
 4.6  Exceções  de  falha  de  página
 A  resposta  do  Xv6  a  exceções  é  bastante  tediosa:  se  uma  exceção  acontece  no  espaço  do  usuário,  o  kernel  encerra  o  processo  
com  falha.  Se  uma  exceção  acontece  no  kernel,  o  kernel  entra  em  pânico.  Operação  real
 48
Machine Translated by Google
 os  sistemas  geralmente  respondem  de  maneiras  muito  mais  interessantes.
 Como  exemplo,  muitos  kernels  usam  falhas  de  página  para  implementar  bifurcações  de  cópia  na  escrita  (COW).  Para  explicar  a  
bifurcação  de  cópia  na  escrita,  considere  a  bifurcação  do  xv6,  descrita  no  Capítulo  3.  A  bifurcação  faz  com  que  o  conteúdo  de  memória  
inicial  do  filho  seja  o  mesmo  do  pai  no  momento  da  bifurcação.  O  Xv6  implementa  a  bifurcação  com  uvmcopy  (kernel/vm.c:306).  que  
aloca  memória  física  para  o  filho  e  copia  a  memória  do  pai  para  ela.  Seria  mais  eficiente  se  o  filho  e  o  pai  pudessem  compartilhar  a  
memória  física  do  pai.  Uma  implementação  direta  disso  não  funcionaria,  no  entanto,  pois  faria  com  que  o  pai  e  o  filho  interrompessem  
a  execução  um  do  outro  com  suas  gravações  na  pilha  e  no  heap  compartilhados.
 Pai  e  filho  podem  compartilhar  memória  física  com  segurança  por  meio  do  uso  apropriado  de  permissões  de  tabela  
de  páginas  e  falhas  de  página.  A  CPU  gera  uma  exceção  de  falha  de  página  quando  um  endereço  virtual  é  usado  sem  
mapeamento  na  tabela  de  páginas,  ou  tem  um  mapeamento  cujo  sinalizador  PTE_V  está  limpo,  ou  um  mapeamento  cujos  
bits  de  permissão  (PTE_R,  PTE_W,  PTE_X,  PTE_U)  proíbem  a  operação  que  está  sendo  tentada.  O  RISC-V  distingue  
três  tipos  de  falha  de  página:  falhas  de  página  de  carregamento  (quando  uma  instrução  de  carregamento  não  consegue  
traduzir  seu  endereço  virtual),  falhas  de  página  de  armazenamento  (quando  uma  instrução  de  armazenamento  não  
consegue  traduzir  seu  endereço  virtual)  e  falhas  de  página  de  instrução  (quando  o  endereço  no  contador  de  programa  
não  traduz).  O  registrador  scause  indica  o  tipo  de  falha  de  página  e  o  registrador  stval  contém  o  endereço  que  não  pôde  
ser  traduzido.
 O  plano  básico  em  uma  bifurcação  COW  é  que  o  pai  e  o  filho  compartilhem  inicialmente  todas  as  páginas  físicas,  mas  cada  um  
mapeie-as  como  somente  leitura  (com  o  sinalizador  PTE_W  limpo).  Pai  e  filho  podem  ler  a  partir  da  memória  física  compartilhada.  Se  
qualquer  um  deles  gravar  uma  determinada  página,  a  CPU  RISC-V  gera  uma  exceção  de  falha  de  página.  O  manipulador  de  traps  do  
kernel  responde  alocando  uma  nova  página  de  memória  física  e  copiando  para  ela  a  página  física  para  a  qual  o  endereço  com  falha  
mapeia.  O  kernel  altera  a  PTE  relevante  na  tabela  de  páginas  do  processo  com  falha  para  apontar  para  a  cópia  e  permitir  gravações,  
bem  como  leituras,  e  então  retoma  o  processo  com  falha  na  instrução  que  causou  a  falha.  Como  a  PTE  permite  gravações,  a  instrução  
reexecutada  agora  será  executada  sem  falhas.  A  cópia  na  gravação  requer  contabilidade  para  ajudar  a  decidir  quando  as  páginas  
físicas  podem  ser  liberadas,  uma  vez  que  cada  página  pode  ser  referenciada  por  um  número  variável  de  tabelas  de  páginas,  
dependendo  do  histórico  de  bifurcações,  falhas  de  página,  execuções  e  saídas.
 Essa  contabilidade  permite  uma  otimização  importante:  se  um  processo  incorrer  em  uma  falha  de  página  de  armazenamento  e  a  
página  física  for  referenciada  apenas  pela  tabela  de  páginas  desse  processo,  nenhuma  cópia  será  necessária.
 A  cópia  na  gravação  torna  a  bifurcação  mais  rápida,  já  que  a  bifurcação  não  precisa  copiar  memória.  Parte  da  memória  precisará  
ser  copiada  posteriormente,  quando  escrita,  mas  frequentemente  a  maior  parte  da  memória  nunca  precisa  ser  copiada.  Um  exemplo  
comum  é  a  bifurcação  seguida  por  um  comando  exec:  algumas  páginas  podem  ser  escritas  após  a  bifurcação,  mas  o  comando  exec  
da  filha  libera  a  maior  parte  da  memória  herdada  da  mãe.
 A  bifurcação  copy-on-write  elimina  a  necessidade  de  copiar  essa  memória.  Além  disso,  a  bifurcação  COW  é  transparente:  nenhuma  
modificação  nos  aplicativos  é  necessária  para  que  eles  se  beneficiem.
 A  combinação  de  tabelas  de  páginas  e  falhas  de  página  abre  uma  ampla  gama  de  possibilidades  interessantes,  além  da  
bifurcação  COW.  Outro  recurso  amplamente  utilizado  é  a  alocação  preguiçosa,  que  possui  duas  partes.  Primeiro,  quando  um  aplicativo  
solicita  mais  memória  chamando  sbrk,  o  kernel  percebe  o  aumento  de  tamanho,  mas  não  aloca  memória  física  e  não  cria  PTEs  para  o  
novo  intervalo  de  endereços  virtuais.  Segundo,  em  uma  falha  de  página  em  um  desses  novos  endereços,  o  kernel  aloca  uma  página  
de  memória  física  e  a  mapeia  na  tabela  de  páginas.  Assim  como  a  bifurcação  COW,  o  kernel  pode  implementar
 49
Machine Translated by Google
 alocação  preguiçosa  de  forma  transparente  para  aplicativos.
 Como  os  aplicativos  frequentemente  solicitam  mais  memória  do  que  precisam,  a  alocação  preguiçosa  é  vantajosa:  o  kernel  não  
precisa  realizar  nenhum  trabalho  para  páginas  que  o  aplicativo  nunca  usa.  Além  disso,  se  o  aplicativo  solicitar  um  aumento  significativo  
no  espaço  de  endereço,  o  sbrk  sem  alocação  preguiçosa  é  caro:  se  um  aplicativo  solicitar  um  gigabyte  de  memória,  o  kernel  precisa  
alocar  e  zerar  262.144  páginas  de  4.096  bytes.  A  alocação  preguiçosa  permite  que  esse  custo  seja  distribuído  ao  longo  do  tempo.  
Por  outro  lado,  a  alocação  preguiçosa  incorre  na  sobrecarga  extra  de  falhas  de  página,  que  envolvem  uma  transição  kernel/usuário.  
Os  sistemas  operacionais  podem  reduzir  esse  custo  alocando  um  lote  de  páginas  consecutivas  por  falha  de  página  em  vez  de  uma  
página  e  especializando  o  código  de  entrada/saída  do  kernel  para  tais  falhas  de  página.
 Outro  recurso  amplamente  utilizado  que  explora  falhas  de  página  é  a  paginação  sob  demanda.  Em  exec,  o  xv6  carrega  todo  o  
texto  e  os  dados  de  um  aplicativo  avidamente  na  memória.  Como  os  aplicativos  podem  ser  grandes  e  a  leitura  do  disco  é  cara,  esse  
custo  inicial  pode  ser  perceptível  para  os  usuários:  quando  o  usuário  inicia  um  aplicativo  grande  a  partir  do  shell,  pode  levar  muito  
tempo  até  que  ele  veja  uma  resposta.  Para  melhorar  o  tempo  de  resposta,  um  kernel  moderno  cria  a  tabela  de  páginas  para  o  espaço  
de  endereço  do  usuário,  mas  marca  as  PTEs  das  páginas  como  inválidas.  Em  uma  falha  de  página,  o  kernel  lê  o  conteúdo  da  página  
do  disco  e  o  mapeia  para  o  espaço  de  endereço  do  usuário.  Assim  como  o  COW  fork  e  a  alocação  preguiçosa,  o  kernel  pode  
implementar  esse  recurso  de  forma  transparente  para  os  aplicativos.
 Os  programas  em  execução  em  um  computador  podem  precisar  de  mais  memória  do  que  a  RAM  disponível  no  computador.
 Para  lidar  com  a  situação  de  forma  elegante,  o  sistema  operacional  pode  implementar  a  paginação  em  disco.  A  ideia  é  armazenar  
apenas  uma  fração  das  páginas  do  usuário  na  RAM  e  armazenar  o  restante  em  disco,  em  uma  área  de  paginação.  O  kernel  marca  
as  PTEs  que  correspondem  à  memória  armazenada  na  área  de  paginação  (e,  portanto,  não  na  RAM)  como  inválidas.  Se  um  aplicativo  
tentar  usar  uma  das  páginas  que  foram  paginadas  para  o  disco,  o  aplicativo  incorrerá  em  uma  falha  de  página  e  a  página  deverá  ser  
paginada:  o  manipulador  de  traps  do  kernel  alocará  uma  página  de  RAM  física,  lerá  a  página  do  disco  para  a  RAM  e  modificará  a  
PTE  relevante  para  apontar  para  a  RAM.
 O  que  acontece  se  uma  página  precisar  ser  paginada,  mas  não  houver  RAM  física  livre?  Nesse  caso,  o  kernel  deve  primeiro  
liberar  uma  página  física,  paginando-a  para  fora  ou  removendo-a  para  a  área  de  paginação  no  disco  e  marcando  as  PTEs  que  se  
referem  a  essa  página  física  como  inválidas.  A  remoção  é  dispendiosa,  portanto,  a  paginação  tem  melhor  desempenho  se  for  pouco  
frequente:  se  os  aplicativos  usarem  apenas  um  subconjunto  de  suas  páginas  de  memória  e  a  união  dos  subconjuntos  couber  na  
RAM.  Essa  propriedade  é  frequentemente  chamada  de  boa  localidade  de  referência.  Assim  como  acontece  com  muitas  técnicas  de  
memória  virtual,  os  kernels  geralmente  implementam  a  paginação  para  o  disco  de  forma  transparente  para  os  aplicativos.
 Os  computadores  frequentemente  operam  com  pouca  ou  nenhuma  memória  física  livre,  independentemente  da  quantidade  de  
RAM  fornecida  pelo  hardware.  Por  exemplo,  provedores  de  nuvem  multiplexam  muitos  clientes  em  uma  única  máquina  para  usar  
seu  hardware  de  forma  econômica.  Outro  exemplo:  usuários  executam  muitos  aplicativos  em  smartphones  com  uma  pequena  
quantidade  de  memória  física.  Em  tais  cenários,  a  alocação  de  uma  página  pode  exigir  a  remoção  prévia  de  uma  página  existente.  
Portanto,  quando  a  memória  física  livre  é  escassa,  a  alocação  é  dispendiosa.
 A  alocação  lenta  e  a  paginação  sob  demanda  são  particularmente  vantajosas  quando  a  memória  livre  é  escassa.
 A  alocação  avida  de  memória  em  sbrk  ou  exec  incorre  no  custo  extra  de  remoção  para  disponibilizar  memória.  Além  disso,  existe  o  
risco  de  que  o  trabalho  avidamente  realizado  seja  desperdiçado,  pois  antes  que  o  aplicativo  use  a  página,  o  sistema  operacional  pode  
tê-la  removido.
 Outros  recursos  que  combinam  exceções  de  paginação  e  falha  de  página  incluem  extensão  automática
 pilhas  e  arquivos  mapeados  na  memória.
 50
Machine Translated by Google
 4.7  Mundo  real
 O  trampolim  e  o  trapframe  podem  parecer  excessivamente  complexos.  Uma  força  motriz  é  que  o  RISC-V  intencionalmente  
faz  o  mínimo  possível  ao  forçar  uma  armadilha,  para  permitir  a  possibilidade  de  um  tratamento  muito  rápido  da  armadilha,  
o  que  acaba  sendo  importante.  Como  resultado,  as  primeiras  instruções  do  manipulador  de  armadilhas  do  kernel  precisam  
ser  executadas  efetivamente  no  ambiente  do  usuário:  a  tabela  de  páginas  do  usuário  e  o  conteúdo  do  registrador  do  
usuário.  E  o  manipulador  de  armadilhas  inicialmente  ignora  fatos  úteis,  como  a  identidade  do  processo  em  execução  ou  o  
endereço  da  tabela  de  páginas  do  kernel.  Uma  solução  é  possível  porque  o  RISC-V  fornece  locais  protegidos  nos  quais  o  
kernel  pode  armazenar  informações  antes  de  entrar  no  espaço  do  usuário:  o  registrador  sscratch  e  as  entradas  da  tabela  
de  páginas  do  usuário  que  apontam  para  a  memória  do  kernel,  mas  são  protegidas  pela  ausência  de  PTE_U.  O  trampolim  
e  o  trapframe  do  Xv6  exploram  esses  recursos  do  RISC-V.
 A  necessidade  de  páginas  especiais  de  trampolim  poderia  ser  eliminada  se  a  memória  do  kernel  fosse  mapeada  na  
tabela  de  páginas  de  usuário  de  cada  processo  (com  os  sinalizadores  de  permissão  PTE  apropriados).  Isso  também  
eliminaria  a  necessidade  de  uma  troca  de  tabela  de  páginas  ao  capturar  do  espaço  do  usuário  para  o  kernel.  Isso,  por  sua  
vez,  permitiria  que  implementações  de  chamadas  de  sistema  no  kernel  aproveitassem  o  mapeamento  da  memória  de  
usuário  do  processo  atual,  permitindo  que  o  código  do  kernel  desreferenciasse  diretamente  os  ponteiros  do  usuário.
 Muitos  sistemas  operacionais  têm  usado  essas  ideias  para  aumentar  a  eficiência.  O  Xv6  as  evita  para  reduzir  as  chances  
de  bugs  de  segurança  no  kernel  devido  ao  uso  inadvertido  de  ponteiros  de  usuário  e  para  reduzir  um  pouco  da  
complexidade  necessária  para  garantir  que  os  endereços  virtuais  do  usuário  e  do  kernel  não  se  sobreponham.
 Sistemas  operacionais  de  produção  implementam  bifurcação  de  cópia  na  gravação,  alocação  preguiçosa,  paginação  
sob  demanda,  paginação  em  disco,  arquivos  mapeados  na  memória,  etc.  Além  disso,  sistemas  operacionais  de  produção  
tentarão  usar  toda  a  memória  física,  seja  para  aplicativos  ou  caches  (por  exemplo,  o  cache  de  buffer  do  sistema  de  
arquivos,  que  abordaremos  mais  adiante  na  Seção  8.2).  O  Xv6  é  ingênuo  nesse  aspecto:  você  quer  que  seu  sistema  
operacional  use  a  memória  física  pela  qual  você  pagou,  mas  o  xv6  não.  Além  disso,  se  o  xv6  ficar  sem  memória,  ele  
retorna  um  erro  para  o  aplicativo  em  execução  ou  o  encerra,  em  vez  de,  por  exemplo,  remover  uma  página  de  outro  
aplicativo
